\section{The Compliance-Credit Mechanism}
\label{sec:credits}

% TODO: policy contribution section — 2-3 pages
%
% The original contribution: market mechanism for adoption incentive
%
% Analogy: EU Emissions Trading System (ETS)
%   - Carbon credits: emit less → receive tradeable credit
%   - Compliance credits: verify more → receive regulatory safe-harbour
%
% Mechanism design:
%   1. Certification: operator deploys BTV stack, audited by
%      accredited third party (analogous to ISO 27001)
%   2. Credit issuance: regulator (DPA, AI Office) issues
%      "Accountability Credit" per verified decision-batch
%   3. Safe-harbour: credit-holding operators receive:
%      (a) reduced enforcement priority
%      (b) presumption of "appropriate technical measures" (GDPR Art.25)
%      (c) reduced AI Act conformity assessment burden (Art. 43)
%   4. Market: credits transferable — large operators can sell
%      excess credits to smaller operators under compliance pressure
%
% Why this doesn't require new law:
%   - GDPR Art.42 already provides for certification schemes
%   - AI Act Art.69 already provides for codes of conduct
%   - Safe-harbour via certification: established precedent (PCI-DSS,
%     SOC 2, ISO 27001 all grant enforcement deference)
%
% Difference from existing proposals (Reisman 2018, Delacroix 2019):
%   - Prior proposals lack the technical substrate to make
%     "certification" meaningful (what exactly is certified?)
%   - BTV stack provides the missing substrate: the certification
%     is of the cryptographic guarantees, not process compliance
%   - This makes the safe-harbour non-gameable: you cannot fake
%     a Merkle inclusion proof
\textbf{[Section under construction]}
